\documentclass[11pt]{article}
\usepackage[T1]{fontenc}
\usepackage{lmodern}
\usepackage[margin=1in]{geometry}
\usepackage{amsmath,amssymb,amsthm,mathtools}
\usepackage{microtype}
\usepackage{needspace}
\usepackage[dvipsnames]{xcolor}
\usepackage{hyperref}
\hypersetup{colorlinks=true,linkcolor=MidnightBlue,citecolor=MidnightBlue,
  urlcolor=MidnightBlue,pdftitle={Scalar compressions and asymptotically optimal tough error models},
  pdfauthor={Yuxuan Zhang},
  pdfsubject={Version 2.0; quantitative construction; internal review}}
\newcommand{\C}{\mathbb C}
\newcommand{\E}{\mathcal E}
\newcommand{\V}{\mathcal V}
\DeclareMathOperator{\Tr}{Tr}
\DeclareMathOperator{\rank}{rank}
\DeclareMathOperator{\ran}{ran}
\DeclareMathOperator{\Span}{span}
\newtheorem{theorem}{Theorem}[section]
\newtheorem{proposition}[theorem]{Proposition}
\newtheorem{lemma}[theorem]{Lemma}
\newtheorem{corollary}[theorem]{Corollary}
\theoremstyle{remark}
\newtheorem{remark}[theorem]{Remark}
\numberwithin{equation}{section}
\setlength{\emergencystretch}{1em}
\title{Scalar compressions and asymptotically optimal\protect\\ tough error models}
\author{Yuxuan Zhang}
\date{16 November 2026 --- Version 2.0\\[4pt]
  \small Internal review; independent proof and priority review pending}

\begin{document}
\maketitle
\begin{abstract}
For an $e$-dimensional complex space of errors acting on $\C^n$, let
$c(e,n)$ be the largest code dimension guaranteed by the Knill--Laflamme
conditions without further assumptions on the errors. We prove
\[
 \max\!\left\{1,\left\lfloor\frac{n+1}{e^2+1}\right\rfloor\right\}
 \le c(e,n)\le
 \left\lceil\frac{n}{e^2+1}\right\rceil
 \qquad (1\le e\le n^2).
\]
The two bounds differ by at most one. Explicit error spaces attain the
upper bound, so their largest codes are within one dimension of the
worst possible value. The lower bound follows from a scalar-compression
theorem: every $D$-dimensional operator system in $M_n(\C)$ has a rank-$k$
scalar compression when $2\le k\le n$ and $n\ge Dk-1$.
We minimize the squared traceless parts of the compressed observables.
At a positive minimum, stationarity supplies two linear dependencies;
these leave a common orthogonal direction along which a sum of second
variations is negative. Thus $c(e,n)=n/(e^2+1)+O(1)$ uniformly in $e,n$,
with relative asymptotic equality whenever $n/(e^2+1)\to\infty$.
We further give a quantitative gradient--curvature inequality for this
variance objective and an explicit approximation iteration. For a normalized
channel with $e\ge2$ specified Kraus operators, $k\ge2$, and
$n\ge e^2k-1$, the iteration constructs a rank-$k$ encoding and a
transpose-type recovery with logical diamond error at most $\varepsilon$
in at most $\lceil768e^2(e^2-1)k^4/\varepsilon^4\rceil$ search steps
in ideal arithmetic. A separate basis-dependent residual guarantee applies
to arbitrary linear error spaces. No finite-precision complexity or exact
finite-step search guarantee is asserted, and some finite integer values
of $c(e,n)$ remain undetermined.
\end{abstract}

\section{The guaranteed code dimension}

How large an exactly correcting subspace must exist if only the dimension
of the error space is known? This is the tough-error-model problem in
Kr\"uger and Werner's collection~\cite[Problem~14]{KW05} and in its
current catalog formulation~\cite{QIQC}.
Let $\E\subseteq M_n(\C)$ be a complex linear space with
$\dim_{\C}\E=e$. A nonzero subspace $C\subseteq\C^n$ corrects $\E$ if
its orthogonal projection $P_C$ satisfies
\begin{equation}\label{eq:kl}
 P_C A^*B P_C=\lambda(A,B)P_C\qquad(A,B\in\E)
\end{equation}
for suitable scalars $\lambda(A,B)$. This is the algebraic
Knill--Laflamme condition~\cite{KL97,KLV00}. Define
\begin{equation}\label{eq:minimax}
 \alpha(\E)=\max\{\dim C:C\text{ satisfies \eqref{eq:kl}}\},
 \qquad
 c(e,n)=\min_{\substack{\E\subseteq M_n(\C)\\\dim_{\C}\E=e}}
 \alpha(\E).
\end{equation}
Every line satisfies \eqref{eq:kl}, so these quantities are positive
integers. The definition imposes neither $I_n\in\E$ nor a
trace-preserving normalization on a basis of $\E$. In particular, $e$
counts a dimension of an arbitrary error space, rather than the Kraus
rank of a normalized channel.

\Needspace{7\baselineskip}
\begin{theorem}\label{thm:main}
For all integers $n\ge1$ and $1\le e\le n^2$,
\begin{equation}\label{eq:main}
 L(e,n):=\max\!\left\{1,\left\lfloor\frac{n+1}{e^2+1}\right\rfloor\right\}
 \le c(e,n)\le
 U(e,n):=\left\lceil\frac{n}{e^2+1}\right\rceil.
\end{equation}
There is an explicit $e$-dimensional error space $\E_{e,n}$ with
$\alpha(\E_{e,n})=U(e,n)$. The gap $U(e,n)-L(e,n)$ is at most one.
\end{theorem}

Thus the construction is asymptotically optimal whenever the guaranteed
code dimension diverges. More precisely,
\begin{equation}\label{eq:asymptotic}
 \left|c(e,n)-\frac{n}{e^2+1}\right|<1,
 \qquad
 \frac{c(e,n)}{n/(e^2+1)}\longrightarrow1
 \quad\text{if }\frac{n}{e^2+1}\longrightarrow\infty.
\end{equation}
The error bound is absolute and allows $e$ to vary with $n$. This answers
the asymptotically matching-bounds version of the problem. It does not
determine the exact rounding at every integer pair.

The new lower bound comes from treating all scalar-compression constraints
at once. For a Hermitian basis $I_n,H_1,\ldots,H_{D-1}$ of an operator
system, consider the sum of the squared traceless parts of
$X^*H_jX$, where $X$ is an $n\times k$ isometry. If this sum is positive
at a minimum, stationarity makes the columns of
$[X,H_1X,\ldots,H_{D-1}X]$ linearly dependent in at least two directions.
When $n\ge Dk-1$, a vector perpendicular to all these columns remains.
Rotating each column of $X$ toward that vector gives a negative sum of
second derivatives, contradicting minimality. Section~\ref{sec:variance}
gives the argument without assumptions on the spectra or commutators of
the observables.

For comparison, the historical estimate recorded in~\cite{KW05}, from
the Knill--Laflamme--Viola construction~\cite{KLV00}, has scale
$n/[e^2(e^2+1)]$. Reimpell also studied this problem through repeated
Hermitian compression and numerical optimization~\cite[Section~3.3.4]{Reimpell}.
In operator-system language, a scalar compression is a quantum anticlique.
Weaver's general existence argument~\cite[Theorem~2.10]{Weaver} uses
diagonal compression followed by a convex partition, with an ambient
dimension requirement quadratic in the operator-system dimension.
Allen and Kornell~\cite[Theorem~2.7]{AK} obtain sharper bounds in small
dimensions: dimension at most $2\ell+1$ guarantees rank $k$ when
$n>3^\ell(k-1)$, and dimension at most $2\ell+2$ when
$n>2\cdot3^\ell(k-1)$. In particular, their result gives
$c(2,n)\ge\lceil n/9\rceil$, whereas \eqref{eq:main} has asymptotic
coefficient $1/5$. These comparisons concern the stated general bounds;
special low-dimensional results can have better integer thresholds than
the theorem below.

The upper construction uses standard rank-one spanning vectors from
Li and Tam~\cite[Section~2.1]{LiTam} and the block-rank obstruction
in Weaver~\cite[Proposition~2.1]{Weaver}. Section~\ref{sec:models}
assembles these ingredients into an $e$-dimensional error space, including
the zero label required by the unnormalized formulation. The construction
and its exact largest-code calculation are included in full. Geometric
and variational descriptions of scalar compression also occur in recent
work on Pauli detection~\cite{Gupta,Sun}; the theorem here is a uniform
existence statement for arbitrary operator systems, with no prescribed
Pauli or representation structure.

Version 2 adds a constructive extension to the existence theorem.
Section~\ref{sec:quantitative} proves a quantitative inequality relating
compression variance to the gradient and the most negative Hessian
curvature. Section~\ref{sec:algorithm} turns it into an explicit ideal-
arithmetic search with a uniform iteration bound, and
Section~\ref{sec:recovery} connects its residual to a physical recovery
error for normalized channels. The variance objective itself is not new:
it appears, for example, in the Pauli-code search of
Gupta et al.~\cite[Section~VI.1]{Gupta}. Nor is transpose recovery new.
The proposed addition is the quantitative analysis at the stated dimension
threshold and its constructive consequence. Comparisons here concern the
specific cited results and do not establish exhaustive priority.

\section{A variational scalar-compression theorem}\label{sec:variance}

An \emph{operator system} $\V\subseteq M_n(\C)$ is a complex linear
subspace containing $I_n$ and satisfying $\V^*=\V$. A rank-$k$
\emph{scalar compression} of $\V$ is an isometry
$X:\C^k\longrightarrow\C^n$ for which $X^*HX$ is a scalar multiple
of $I_k$ for every $H\in\V$. Equivalently, $P=XX^*$ is a rank-$k$
projection with $P\V P=\C P$.

\Needspace{10\baselineskip}
\begin{theorem}\label{thm:compression}
Let $\V\subseteq M_n(\C)$ be an operator system of complex dimension
$D$. If $2\le k\le n$ and
\begin{equation}\label{eq:threshold}
 n\ge Dk-1,
\end{equation}
then $\V$ has a rank-$k$ scalar compression. Consequently, every such
system has a scalar compression of rank at least
\begin{equation}\label{eq:rank-bound}
 \min\!\left\{n,\max\!\left\{1,
 \left\lfloor\frac{n+1}{D}\right\rfloor\right\}\right\}.
\end{equation}
For $D=1$, the entire space is a scalar compression.
\end{theorem}

\begin{proof}
The self-adjoint part of $\V$ has real dimension $D$ and spans $\V$
over $\C$. Choose a Hermitian basis $I_n,H_1,\ldots,H_{D-1}$.
For an isometry $X\in M_{n\times k}(\C)$, set
\begin{equation}\label{eq:objective}
 G_j=X^*H_jX,\qquad
 \mu_j=\frac{\Tr G_j}{k},\qquad
 T_j=G_j-\mu_j I_k,\qquad
 F(X)=\sum_{j=1}^{D-1}\Tr(T_j^2).
\end{equation}
Each $T_j$ is Hermitian and traceless. Thus $F(X)\ge0$, with equality
precisely when $X$ is a scalar compression. The function is continuous
on the compact space of isometries and invariant under $X\mapsto XU$
for unitary $U\in M_k(\C)$. Fix a global minimizer $X$ and suppose
that $F(X)>0$.

For an arbitrary variation $Z$ of $X$, differentiation gives
\begin{equation}\label{eq:differential}
 dF_X[Z]=4\operatorname{Re}\Tr\!\left[
 \left(\sum_jH_jXT_j\right)^*Z\right].
\end{equation}
Every $Z$ satisfying $X^*Z=0$ is tangent to the isometries at $X$.
Since \eqref{eq:differential} vanishes for all such $Z$, the horizontal
component of the gradient is zero. Hence
\begin{equation}\label{eq:stationarity}
 \sum_j H_jXT_j=XS,\qquad
 S=\sum_jG_jT_j.
\end{equation}
Each $G_j$ commutes with $T_j$, so $S$ is Hermitian. Introduce the block
matrices
\begin{equation}\label{eq:dependencies}
 M=\begin{bmatrix}X&H_1X&\cdots&H_{D-1}X\end{bmatrix},
 \qquad
 B=\begin{bmatrix}-S\\T_1\\\vdots\\T_{D-1}\end{bmatrix}.
\end{equation}
Equation~\eqref{eq:stationarity} says $MB=0$. Since $F(X)>0$, at least
one $T_j$ is a nonzero traceless Hermitian matrix. Such a matrix has
both a positive and a negative eigenvalue, and therefore rank at least
two. As a block of $B$, it implies $\rank B\ge2$. The matrix $M$ has
$Dk$ columns, and its kernel contains the range of $B$. Consequently,
\begin{equation}\label{eq:rank-defect}
 \rank M\le Dk-2<n.
\end{equation}
Choose a unit vector $y\in\C^n$ orthogonal to the range of $M$. Then
\begin{equation}\label{eq:orthogonality}
 X^*y=0,\qquad X^*H_jy=0\quad(1\le j\le D-1),
\end{equation}
where the second equality uses $H_j^*=H_j$.

Write $X=[x_1,\ldots,x_k]$. For each $a\in\{1,\ldots,k\}$, replace
the column $x_a$ by $\cos\theta\,x_a+\sin\theta\,y$, leaving all
other columns fixed. Denote this isometry by $X_a(\theta)$ and put
$F_a(\theta)=F(X_a(\theta))$. Let $f_a$ be the $a$th coordinate vector
of $\C^k$, $E_a=f_af_a^*$, and
$C_a(\theta)=I_k+(\cos\theta-1)E_a$. Then
\[
 X_a(\theta)=XC_a(\theta)+\sin\theta\,yf_a^*.
\]
Writing $h_j=y^*H_jy$, the orthogonality in
\eqref{eq:orthogonality} gives the exact identity
\begin{equation}\label{eq:rotation-exact}
 X_a(\theta)^*H_jX_a(\theta)
 =C_a(\theta)G_jC_a(\theta)+\sin^2\theta\,h_jE_a.
\end{equation}
Its expansion at zero is
\begin{equation}\label{eq:rotation-expansion}
 G_j+\theta^2\left(h_jE_a-\tfrac12(E_aG_j+G_jE_a)\right)
 +O(\theta^4).
\end{equation}
There is no first-order term. Since
$\Tr(T_j^2)=\Tr(G_j^2)-(\Tr G_j)^2/k$, differentiation yields
\begin{align}
 F_a''(0)
 &=4\sum_j\Tr\!\left[T_j
 \left(h_jE_a-\tfrac12(E_aG_j+G_jE_a)\right)\right]\notag\\
 &=4\sum_j\left(h_j(T_j)_{aa}-(G_jT_j)_{aa}\right).
 \label{eq:second-variation}
\end{align}
The last step again uses $G_jT_j=T_jG_j$. Summing over the $k$ column
rotations cancels the terms containing $h_j$, because $\Tr T_j=0$.
Moreover, $\Tr(G_jT_j)=\Tr(T_j^2)$. Therefore
\begin{equation}\label{eq:negative-sum}
 \sum_{a=1}^kF_a''(0)
 =-4\sum_j\Tr(T_j^2)=-4F(X)<0.
\end{equation}
Every $F_a''(0)$ must be nonnegative at a minimum, which is a
contradiction. Thus $F(X)=0$.

A rank-one scalar compression always exists. If $D\ge2$ and the rank
in \eqref{eq:rank-bound} is at least two, it satisfies
\eqref{eq:threshold} and is at most $n$, so the first part applies.
If $D=1$, then $\V=\C I_n$ and rank $n$ is immediate.
\end{proof}

\begin{remark}
The same argument shows that any positive stationary value of $F$ has
a direction of negative second variation under \eqref{eq:threshold}.
The existence proof uses compactness to select a global minimum.
Sections~\ref{sec:quantitative}--\ref{sec:recovery} supply a quantitative
approximation search, while exact finite-step construction remains open.
\end{remark}

\begin{corollary}\label{cor:lower}
Every $e$-dimensional complex error space in $M_n(\C)$ has a correcting
code of dimension at least $L(e,n)$ in \eqref{eq:main}.
\end{corollary}

\begin{proof}
Choose a complex basis $A_1,\ldots,A_e$ of $\E$ and form
\begin{equation}\label{eq:error-system}
 \V_{\E}=\Span_{\C}\{I_n,A_i^*A_j:1\le i,j\le e\}.
\end{equation}
This is an operator system of dimension $d\le e^2+1$.
Let $k=\lfloor(n+1)/(e^2+1)\rfloor$. If $k\ge2$, then
$k\le n$ and $n\ge(e^2+1)k-1\ge dk-1$, so
Theorem~\ref{thm:compression} gives a rank-$k$ scalar compression of
$\V_{\E}$. The pairwise products $A_i^*A_j$ then satisfy
\eqref{eq:kl}, and sesquilinearity extends the condition to all
$A,B\in\E$. When $k\le1$, any line is a correcting code.
\end{proof}

\section{Explicit tough error spaces}\label{sec:models}

We next construct the error spaces attaining the upper bound in
Theorem~\ref{thm:main}. The basic obstruction is the following
block-rank calculation~\cite[Proposition~2.1]{Weaver}.

\Needspace{7\baselineskip}
\begin{proposition}\label{prop:blocks}
Let $P_1,\ldots,P_t$ be nonzero mutually orthogonal projections summing
to $I_n$, and let $\V=\Span_{\C}\{P_1,\ldots,P_t\}$. The largest
rank of a scalar compression of $\V$ is
$\max_s\rank P_s$.
\end{proposition}

\begin{proof}
If a rank-$k$ projection $P$ satisfies $PP_sP=\lambda_sP$ for all $s$,
positivity gives $\lambda_s\ge0$ and summing gives $\sum_s\lambda_s=1$.
For some $s$, $\lambda_s>0$, and hence
$k=\rank(PP_sP)\le\rank P_s$. Conversely, the range of any $P_s$
is a scalar compression of $\V$.
\end{proof}

Assume first that $1\le e<n$, and put $t=\min\{n,e^2+1\}$.
Let $u_1,\ldots,u_e$ be the coordinate vectors in $\C^e$. Make an
ordered list of $e^2$ vectors by taking all $u_i$ first, followed, in
lexicographic order on $i<j$, by
\begin{equation}\label{eq:labels}
 u_i+u_j,\qquad u_i+\mathrm{i}u_j.
\end{equation}
These are the coordinate and pair vectors in the rank-one spanning
construction of Li and Tam~\cite[Section~2.1]{LiTam}, without normalizing
their lengths. Let $v_0=0$, and let $v_1,\ldots,v_{t-1}$ be the first
$t-1$ vectors of the list. Because $t-1\ge e$, all coordinate vectors
are included.

Decompose $\C^n$ into $t$ nonzero coordinate blocks with dimensions
differing by at most one, and denote their projections by
$P_0,\ldots,P_{t-1}$. Explicitly, writing $n=qt+r$ with $0\le r<t$,
take $r$ blocks of dimension $q+1$ and $t-r$ blocks of dimension $q$.
Define diagonal errors
\begin{equation}\label{eq:errors}
 A_a=\sum_{s=0}^{t-1}(v_s)_aP_s\quad(1\le a\le e),
 \qquad \E_{e,n}=\Span_{\C}\{A_1,\ldots,A_e\}.
\end{equation}
On the block labelled by $u_i$, $A_a$ equals $\delta_{ai}I$.
Thus the $A_a$ are linearly independent and $\dim\E_{e,n}=e$.

To identify their product space, write $D_i$ for the projection of the
block labelled by $u_i$. For $i<j$, let $R_{ij}$ and $S_{ij}$ be the
projections of the blocks labelled by $u_i+u_j$ and
$u_i+\mathrm{i}u_j$, respectively; a missing block means the zero
projection. The labels give
\begin{equation}\label{eq:cross-products}
 A_i^*A_j=R_{ij}+\mathrm{i}S_{ij},\qquad
 R_{ij}=\frac{A_i^*A_j+A_j^*A_i}{2},\qquad
 S_{ij}=\frac{A_i^*A_j-A_j^*A_i}{2\mathrm{i}}.
\end{equation}
They also give
\begin{align}
 D_i&=A_i^*A_i-
 \sum_{\substack{p<q\\i\in\{p,q\}}}(R_{pq}+S_{pq}),
 \label{eq:coordinate-blocks}\\
 P_0&=I_n-\sum_iD_i-\sum_{i<j}(R_{ij}+S_{ij}).
 \label{eq:zero-block}
\end{align}
Every block projection therefore belongs to $\V_{\E_{e,n}}$.
Conversely, the errors and all their products are constant on each
block. We have proved
\begin{equation}\label{eq:block-system}
 \V_{\E_{e,n}}=\Span_{\C}\{P_0,\ldots,P_{t-1}\}.
\end{equation}
Proposition~\ref{prop:blocks} now gives
\begin{equation}\label{eq:largest-code}
 \alpha(\E_{e,n})=\left\lceil\frac nt\right\rceil
 =\left\lceil\frac{n}{e^2+1}\right\rceil.
\end{equation}
The last equality holds because $t=e^2+1$ when $n\ge e^2+1$, and
both ceilings equal one otherwise. The rank argument applies to every
subspace, including codes with support across several blocks.

For $e\ge n$, let $\E_{e,n}$ contain all $n$ diagonal matrix units
and extend them by off-diagonal matrix units, in lexicographic order,
to a linearly independent set of $e$ errors. Since $\V_{\E_{e,n}}$
contains the $n$ rank-one coordinate projections, the same rank argument
forces every code to have dimension one. Here
$\lceil n/(e^2+1)\rceil=1$. This covers the remaining parameters,
including $n=e=1$, and proves the upper bound in
Theorem~\ref{thm:main}.

The smallest two-error example displaying all five labels is
\begin{equation}\label{eq:five-block-example}
 A_1=\operatorname{diag}(0,1,0,1,1),\qquad
 A_2=\operatorname{diag}(0,0,1,1,\mathrm{i}).
\end{equation}
Its product system, with the identity adjoined as in
\eqref{eq:error-system}, is the full diagonal algebra in $M_5(\C)$,
so its largest code has dimension one. Replacing the five coordinates
by balanced blocks gives largest code dimension $\lceil n/5\rceil$
for every $n\ge5$.

\section{A quantitative scalar-compression inequality}\label{sec:quantitative}

We now strengthen the existence argument to a uniform estimate away from
stationarity. Let $H_1,\ldots,H_m$ be Hermitian $n\times n$ matrices;
they need not be linearly independent. Throughout this section assume
$m\ge1$, $2\le k\le n$, and
\begin{equation}\label{eq:quantitative-dimension}
 n\ge(m+1)k-1.
\end{equation}
Use the definitions $G_j,T_j,F$ in \eqref{eq:objective}, now summing
over $1\le j\le m$. The Grassmannian is represented by isometries $X$,
with horizontal tangent space $X^*Z=0$ and real metric
$\langle Z,W\rangle=\operatorname{Re}\Tr(Z^*W)$. Thus
\begin{equation}\label{eq:residual-gradient}
 R=(I_n-XX^*)\sum_jH_jXT_j,\qquad \operatorname{grad}F=4R.
\end{equation}
The Hessian quadratic form at $Z$ is the second derivative along the
geodesic $\exp[t(ZX^*-XZ^*)]X$. Define
\begin{equation}\label{eq:quantitative-constants}
 \rho=\|R\|_{\rm op},\qquad
 \eta=\max\{0,-\lambda_{\min}(\operatorname{Hess}F)\},\qquad
 a=(k-1)\sqrt{k+1}.
\end{equation}
All Hessian eigenvalues refer to this real metric.

\begin{theorem}\label{thm:quantitative}
Under \eqref{eq:quantitative-dimension}, every isometry $X$ satisfies
\begin{align}
 F^2&\le a^2\rho^2+\frac{k}{4}\eta F,\label{eq:robust-quadratic}\\
 F&\le a\rho+\frac{k}{4}\eta
 \le\frac{a}{4}\|\operatorname{grad}F\|_{\rm F}
       +\frac{k}{4}\eta.\label{eq:robust-linear}
\end{align}
In particular, if $F>0$ at a stationary point, then
$\lambda_{\min}(\operatorname{Hess}F)\le-4F/k$.
\end{theorem}

\begin{lemma}\label{lem:stack}
For traceless Hermitian $k\times k$ matrices $T_j$, let
$A=[T_1;\ldots;T_m]$ and $F=\sum_j\|T_j\|_{\rm F}^2$.
The second largest singular value satisfies
\begin{equation}\label{eq:stack-bound}
 \sigma_2(A)^2\ge\frac{F}{k(k-1)}.
\end{equation}
The constant is sharp for arbitrary such stacks.
\end{lemma}

\begin{proof}
Given a unit vector $v$, write a traceless Hermitian $T$ in a basis
beginning with $v$ as $\left[\begin{smallmatrix}b&w^*\\w&C\end{smallmatrix}\right]$.
Here $b\in\mathbb R$ and $\Tr C=-b$, so
\[
 \|T\|_{\rm F}^2
 =b^2+2\|w\|^2+\|C\|_{\rm F}^2
 \ge\frac{k}{k-1}b^2+2\|w\|^2
 \ge\frac{k}{k-1}\|Tv\|^2.
\]
For $Q=A^*A=\sum_jT_j^2$, this gives
$\lambda_{\max}(Q)\le(k-1)F/k$. Since $\Tr Q=F$, the remaining
$k-1$ eigenvalues sum to at least $F/k$. Their largest is at least
$F/[k(k-1)]$, proving \eqref{eq:stack-bound}. The single matrix
$\operatorname{diag}(k-1,-1,\ldots,-1)$ attains equality.
\end{proof}

Complete $X$ to a unitary $[X,Y]$, and write
\begin{equation}\label{eq:block-hessian-data}
 [X,Y]^*H_j[X,Y]=
 \begin{bmatrix}G_j&C_j^*\\C_j&D_j\end{bmatrix},\quad
 C=[C_1,\ldots,C_m],\quad S=\sum_jG_jT_j.
\end{equation}
Here $C$ has $q=n-k$ rows and $p=mk$ columns; \eqref{eq:quantitative-dimension}
says $q\ge p-1$. The identity $CA=Y^*R$ supplies a useful
approximately common orthogonal direction whenever $F>0$.
Indeed, take the two leading left singular vectors of $A$ as the
columns of $U_2$. Then
$\|CU_2\|_{\rm op}\le\rho/\sigma_2(A)$, by restricting
$CA$ to the corresponding two right singular vectors.
The rank of $C(I_p-U_2U_2^*)$ is at most $p-2<q$, so it has
a unit vector $u$ in its left kernel. Consequently
\begin{equation}\label{eq:approx-common}
 \|C^*u\|^2\le\|CU_2\|_{\rm op}^2
 \le\frac{k(k-1)\rho^2}{F}.
\end{equation}
A unit vector minimizing $\|C^*u\|$ satisfies the same bound.
This includes rectangular $C$ with $q=p-1$; no inverse of $A$ is used.

For completeness, the full Hessian calculation needed here does not
assume stationarity or common orthogonality. Write a horizontal direction
as $Z=YW$ and put $L_j=C_j^*W+W^*C_j$, with
$L_{j,0}=L_j-(\Tr L_j)I_k/k$. Along its geodesic,
$X'(0)=YW$ and $X''(0)=-XW^*W$. Therefore
\[
 G_j'(0)=L_j,\qquad
 G_j''(0)=2W^*D_jW-W^*WG_j-G_jW^*W.
\]
Differentiating $F=\sum_j\Tr(T_j^2)$ gives
\begin{equation}\label{eq:full-hessian}
 \operatorname{Hess}F[YW,YW]
 =2\sum_j\|L_{j,0}\|_{\rm F}^2
  +4\sum_j\Tr(T_jW^*D_jW)-4\Tr(SW^*W).
\end{equation}
Let $y=Yu$ for a unit $u$ and let $f_b$ denote the $b$th coordinate
vector in $\C^k$. Apply \eqref{eq:full-hessian} to the $2k$ unit
directions $Z_b=yf_b^*$ and $\mathrm{i}Z_b$. For $z=C_j^*u$, the
two first variations are $zf_b^*+f_bz^*$ and
$\mathrm{i}(zf_b^*-f_bz^*)$. Summing their squared traceless norms
over $b$ gives $4(k-1/k)\|z\|^2$. The $D_j$ terms cancel since
$\Tr T_j=0$, and the last terms sum to $-8\Tr S=-8F$. Hence
\begin{equation}\label{eq:phase-sum}
 \sum_{b=1}^k\bigl(\operatorname{Hess}F[Z_b,Z_b]
             +\operatorname{Hess}F[\mathrm{i}Z_b,\mathrm{i}Z_b]\bigr)
 =8\left(k-\frac1k\right)\|C^*u\|^2-8F.
\end{equation}
Both phases are needed for this identity at a general point.

\begin{proof}[Proof of Theorem~\ref{thm:quantitative}]
The case $F=0$ is immediate. Otherwise choose $u$ as in
\eqref{eq:approx-common}. The left side of \eqref{eq:phase-sum}
is at least $-2k\eta$, while its right side is at most
$8a^2\rho^2/F-8F$. Rearrangement proves
\eqref{eq:robust-quadratic}. Its positive root is at most
$a\rho+k\eta/4$; also
$\rho\le\|\operatorname{grad}F\|_{\rm F}/4$.
At stationarity $\rho=0$, so a positive $F$ forces $\eta\ge4F/k$.
\end{proof}

If $n\ge(m+1)k$, then $q\ge p$ and a literal left kernel is
available when $q>p$. At $q=p$ it need not be. The uniform argument
above covers both cases; it requires no rank assumption on $C$.
If $q>p$, \eqref{eq:phase-sum} directly gives a curvature at most
$-4F/k$ at every point with $F>0$, even without stationarity.

\section{An explicit approximation iteration}\label{sec:algorithm}

Fix bounds $h_j\ge\|H_j\|_{\rm op}$ and let
\begin{equation}\label{eq:beta}
 \beta=k\sum_{j=1}^m h_j^2.
\end{equation}
Then $0\le F\le\beta$. The following construction counts ideal
real and complex matrix operations, including singular vectors and
geodesic evaluations. It is not a finite-precision complexity claim.

\begin{theorem}\label{thm:iteration}
Assume \eqref{eq:quantitative-dimension}. For every $\tau>0$, the
iteration below starts from any isometry and attains $F\le\tau$
in at most
\begin{equation}\label{eq:iteration-count}
 \left\lceil\frac{48\beta k^3}{\tau}\right\rceil
\end{equation}
iterations in ideal arithmetic. If $m=0$ or $\beta=0$, $F=0$ and
no iteration is needed.
\end{theorem}

At a current isometry $X$, stop if $F\le\tau$; otherwise compute
$g=\|\operatorname{grad}F\|_{\rm F}$ and $a$ from
\eqref{eq:quantitative-constants}.
\begin{enumerate}
\item If $g\ge F/a$, use the unit direction
$Z=-\operatorname{grad}F/g$ and replace $X$ by
$\exp[t(ZX^*-XZ^*)]X$, with $t=g/(16\beta)$.
\item If $g<F/a$, form $C$ in \eqref{eq:block-hessian-data} and
find a unit vector minimizing $\|C^*u\|$. Among the $2k$ directions
in \eqref{eq:phase-sum}, choose one with the smallest curvature $q_0$.
Evaluate its geodesic at both $t$ and $-t$, where
$t^2=F/(32\beta k)$, and retain the point with smaller $F$.
\end{enumerate}
The second step needs a singular-vector calculation and $2k$ directional
curvatures; it does not require diagonalizing the full Hessian.

\begin{proof}
For any unit horizontal $Z$, the skew-Hermitian generator
$J=ZX^*-XZ^*$ satisfies $\|J\|_{\rm op}=\|Z\|_{\rm op}\le1$.
Along $X(t)=\exp(tJ)X$, each derivative $J^rX(t)$ has operator norm
at most one and Frobenius norm at most $\sqrt{k}$. Leibniz' rule yields
\[
 \left\|\frac{d^r}{dt^r}\bigl(X(t)^*H_jX(t)\bigr)\right\|_{\rm F}
 \le 2^r\sqrt{k}\,h_j.
\]
Traceless projection is a Frobenius contraction. Differentiating its
squared norm once more gives, for all $r\ge0$ and all $t$,
\begin{equation}\label{eq:derivative-bound}
 \left|\frac{d^r}{dt^r}F(X(t))\right|\le4^r\beta.
\end{equation}
In the first branch, $F'(0)=-g$ and $|F''(t)|\le16\beta$, so
\begin{equation}\label{eq:gradient-decrease}
 F(X(t))\le F-\frac{g^2}{32\beta}
 \le F-\frac{F^2}{32\beta a^2}
 \le F-\frac{F^2}{32\beta k^3}.
\end{equation}
In the second branch, $\rho\le g/4<F/(4a)$. Combining
\eqref{eq:approx-common} and \eqref{eq:phase-sum}, and dividing by
the $2k$ directions, gives
\begin{equation}\label{eq:chosen-curvature}
 q_0<-\frac{15F}{4k}\le-\frac{2F}{k}.
\end{equation}
Average the two Taylor expansions at $t$ and $-t$; their odd terms
cancel. Using $|F^{(4)}|\le256\beta$ from
\eqref{eq:derivative-bound} gives
\begin{align}
 \frac{F(X(t))+F(X(-t))}{2}
 &\le F+\frac{q_0t^2}{2}+\frac{256\beta}{24}t^4\notag\\
 &\le F-\frac{F^2}{48\beta k^2}.\label{eq:curvature-decrease}
\end{align}
One of the two tested values is no greater than their average. Both
branches therefore satisfy
\begin{equation}\label{eq:uniform-decrease}
 F_{\rm next}\le F-\frac{F^2}{48\beta k^3}.
\end{equation}
Unless $F_{\rm next}=0$, this implies
$1/F_{\rm next}\ge1/F+1/(48\beta k^3)$; here $F\le\beta$ ensures
the denominator is positive. Iterating proves \eqref{eq:iteration-count}.
\end{proof}

\subsection{Residuals for an arbitrary linear error space}

Fix a basis $A_1,\ldots,A_e$ for the original, unnormalized error space.
Use $m=e^2$ Hermitian observables
\begin{equation}\label{eq:arbitrary-observables}
 A_i^*A_i\ (1\le i\le e),\qquad
 \frac{A_i^*A_j+A_j^*A_i}{\sqrt2},\quad
 \frac{A_i^*A_j-A_j^*A_i}{\mathrm{i}\sqrt2}\quad(i<j).
\end{equation}
Their objective is exactly
\begin{equation}\label{eq:kl-residual}
 F=\sum_{i,j=1}^e
 \|X^*A_i^*A_jX-\lambda_{ij}I_k\|_{\rm F}^2,\qquad
 \lambda_{ij}=\frac{\Tr(X^*A_i^*A_jX)}{k}.
\end{equation}
To see this, pair the $(i,j)$ and $(j,i)$ residuals, which are adjoints,
and use the Frobenius parallelogram identity for their real and imaginary
Hermitian parts. Thus for $n\ge(e^2+1)k-1$ the iteration produces
$F\le\delta^2$ for any $\delta>0$. If $A=\sum_i a_iA_i$ and
$B=\sum_i b_iA_i$, Cauchy--Schwarz gives
\begin{equation}\label{eq:arbitrary-kl-bound}
 \|X^*A^*BX-(a^*\lambda b)I_k\|_{\rm op}
 \le\delta\|a\|_2\|b\|_2.
\end{equation}
The coefficient norms and $\beta$ refer to the specified basis.
Rescaling that basis changes them. In particular, this algebraic
statement makes no assertion about a physical channel's recovery error
without a channel normalization.

\section{A recovery channel with an explicit error bound}\label{sec:recovery}

We next assume that $A_1,\ldots,A_e\in M_n(\C)$ are Kraus operators
with $\sum_iA_i^*A_i=I_n$, and write
$\mathcal N(\sigma)=\sum_iA_i\sigma A_i^*$. They need not be
linearly independent; here $e$ counts the specified Kraus operators.
For an isometry $X:\C^k\to\C^n$, let
$\mathcal E_X(\rho)=X\rho X^*$ and set
\begin{equation}\label{eq:gram-data}
 \begin{split}
 W&=[A_1X,\ldots,A_eX],\qquad K=W^*W,\\
 \Lambda_{ij}&=\frac{\Tr(K_{ij})}{k},\qquad
 L=\Lambda\otimes I_k,\qquad F=\|K-L\|_{\rm F}^2.
 \end{split}
\end{equation}
Thus $F$ equals \eqref{eq:kl-residual}, and
$\Lambda\ge0$, $\Tr\Lambda=1$.

\begin{theorem}\label{thm:recovery}
Assume $n\ge ek$. Extend the polar factor of $W$ to an isometry
$U:\C^e\otimes\C^k\to\C^n$, so that $W=U\sqrt K$.
For any state $\rho_*$ on $\C^k$, the map
\begin{equation}\label{eq:decoder}
 \mathcal D(\sigma)=\Tr_{\C^e}(U^*\sigma U)
       +\Tr[(I_n-UU^*)\sigma]\rho_*
\end{equation}
is a channel on all of $M_n(\C)$ and satisfies
\begin{equation}\label{eq:diamond-bound}
 \|\mathcal D\mathcal N\mathcal E_X-\operatorname{id}_k\|_\diamond
 \le\min\{2,\,2\sqrt e\,F^{1/4}\}.
\end{equation}
The diamond norm uses the full trace-norm convention, with no factor
of $1/2$. No invertibility of $K$ or $\Lambda$ is required.
\end{theorem}

\begin{proof}
The matrix $\Lambda$ is positive because it is the normalized partial
trace of the positive block matrix $K$; channel normalization gives
$\Tr\Lambda=1$. The partial polar isometry of $W$ maps
$\operatorname{supp}K$ into $\C^n$. The inequality $n\ge ek$
allows it to be completed isometrically on the kernel of $K$.
Both terms of \eqref{eq:decoder} are completely positive: the second
is a positive measurement followed by preparation of $\rho_*$. Their
traces sum to $\Tr\sigma$, so $\mathcal D$ is trace preserving.

We use the square-root estimate
\begin{equation}\label{eq:sqrt-holder}
 \|\sqrt K-\sqrt L\|_{\rm op}
 \le\|K-L\|_{\rm op}^{1/2}
 \qquad(K,L\ge0).
\end{equation}
Here is an argument that also covers singular matrices. If
$d=\|K-L\|_{\rm op}$, then $K\le L+dI$ and $L\le K+dI$.
The square root is operator monotone, as follows from the convergent
integral
\[
 \sqrt A=\frac1\pi\int_0^\infty
 t^{-1/2}A(A+tI)^{-1}\,dt
\]
and $A(A+tI)^{-1}=I-t(A+tI)^{-1}$. Also
$\sqrt{L+dI}\le\sqrt L+\sqrt d\,I$ by scalar functional calculus
of $L$. Applying these facts in both orders bounds both signs of
$\sqrt K-\sqrt L$ by $\sqrt d\,I$, proving \eqref{eq:sqrt-holder}.

Put $W_0=U\sqrt L$, and split it into $e$ blocks $B_i$ of size
$n\times k$. Explicitly,
$B_i=U((\sqrt\Lambda\,f_i)\otimes I_k)$, so
\begin{equation}\label{eq:nearby-channel}
 \mathcal N_0(\rho)=\sum_iB_i\rho B_i^*
 =U(\Lambda\otimes\rho)U^*,\qquad
 \mathcal D\mathcal N_0=\operatorname{id}_k.
\end{equation}
The map $\mathcal N_0$ is trace preserving since $\Tr\Lambda=1$.
Equation~\eqref{eq:sqrt-holder} implies
$\|W-W_0\|_{\rm op}\le F^{1/4}$.
For $\Delta_i=A_iX-B_i$, choose the Stinespring isometries
$V\psi=\sum_i A_iX\psi\otimes f_i$ and
$V_0\psi=\sum_iB_i\psi\otimes f_i$. Then
\begin{equation}\label{eq:stinespring-distance}
 \|V-V_0\|_{\rm op}^2
 =\left\|\sum_i\Delta_i^*\Delta_i\right\|_{\rm op}
 \le e\|W-W_0\|_{\rm op}^2\le e\sqrt F.
\end{equation}
For an arbitrary reference and trace-class input $Z$, expand
$VZV^*-V_0ZV_0^*$ as
$(V-V_0)ZV^*+V_0Z(V^*-V_0^*)$, with reference identities understood.
The trace-norm ideal property and contraction under the partial trace
give
$\|\mathcal N\mathcal E_X-\mathcal N_0\|_\diamond
 \le2\|V-V_0\|_{\rm op}$.
Postcomposition by $\mathcal D$ is also contractive. Combining this
with \eqref{eq:nearby-channel} and \eqref{eq:stinespring-distance}
proves the nontrivial term in \eqref{eq:diamond-bound}; two channels
have diamond distance at most two.
\end{proof}

The recovery construction is an established one. On the range of $W$,
the polar factor is $(WW^*)^{-1/2}W$, with the inverse restricted to
the support. The logical recovery Kraus operators are consequently
\[
 X^*A_i^*\,[\mathcal N(XX^*)]^{-1/2}.
\]
These are the transpose/Petz recovery operators for the maximally mixed
code state~\cite[Eq.~(4)]{NgMandayam}. Completion of $U$ only specifies
the action outside the actually noisy code support. General approximate
Knill--Laflamme conditions and near-optimal recovery are treated
in~\cite{BenyOreshkov}; Stinespring continuity is developed
in~\cite{KSW}. We use an elementary, conservative bound expressed in
the particular variance residual controlled by our iteration, without
claiming a new general recovery principle or an optimal constant.

\begin{corollary}\label{cor:constructive-channel}
Let $e\ge2$, $k\ge2$, $n\ge e^2k-1$, and
$\sum_{i=1}^e A_i^*A_i=I_n$. For every $0<\varepsilon\le2$, an
isometry $X$ and a recovery channel $\mathcal D$ with logical diamond
error at most $\varepsilon$ can be constructed using at most
\begin{equation}\label{eq:physical-iteration-count}
 \left\lceil\frac{768e^2(e^2-1)k^4}{\varepsilon^4}\right\rceil
\end{equation}
search iterations in ideal arithmetic, followed by a polar decomposition
and isometric completion.
\end{corollary}

\begin{proof}
Choose a Hilbert--Schmidt orthonormal Hermitian basis
$B_0,\ldots,B_{e^2-1}$ of $M_e(\C)$ with $B_0=I_e/\sqrt e$, and set
\begin{equation}\label{eq:normalized-observables}
 H_\alpha=\sum_{i,j}(B_\alpha)_{ij}A_i^*A_j.
\end{equation}
The map $B\mapsto\sum_{i,j}B_{ij}A_i^*A_j$ is unital and completely
positive: with $V_A\psi=\sum_i A_i\psi\otimes f_i$, it is
$V_A^*(I_n\otimes B)V_A$. Thus each $H_\alpha$ is Hermitian and
$\|H_\alpha\|_{\rm op}\le\|B_\alpha\|_{\rm op}\le1$.
Also $H_0=I_n/\sqrt e$, whose traceless compression vanishes.
The coefficient transformation in \eqref{eq:normalized-observables}
is unitary on the $e^2$ matrix entries, so it preserves the sum of
squared Frobenius residuals. The remaining $m=e^2-1$ observables
therefore have exactly the objective $F$ in \eqref{eq:gram-data}.

Apply Theorem~\ref{thm:iteration} with
$\beta=k(e^2-1)$ and $\tau=\varepsilon^4/(16e^2)$.
Its dimension condition is precisely $n\ge e^2k-1$, and its count
is \eqref{eq:physical-iteration-count}. Since $e\ge2$ and $k\ge2$,
this dimension condition also implies $n\ge ek$. The recovery in
Theorem~\ref{thm:recovery} now has error at most $\varepsilon$.
\end{proof}

For one square Kraus operator, normalization makes the channel unitary,
so its inverse corrects every code. For $k=1$, a fixed logical state
can be recovered directly. The nontrivial search statement above counts
operations on dense matrices, not operations polynomial in the number
of qubits. It does not establish a finite-precision bit bound or exact
finite termination at $\varepsilon=0$. The arbitrary-error result in
\eqref{eq:arbitrary-kl-bound} retains the threshold
$n\ge(e^2+1)k-1$; the normalized threshold in this corollary cannot be
substituted into the unnormalized minimax problem.

\section{Asymptotics and the remaining integer gap}\label{sec:integer-gap}

It remains to check the one-dimension gap in Theorem~\ref{thm:main}.
Write $D=e^2+1$ and $n=qD+r$, with $0\le r<D$.
If $r=0$ or $r=D-1$, the two sides of \eqref{eq:main} agree.
For $1\le r\le D-2$ and $q\ge1$, they are $q$ and $q+1$.
If $q=0$, both bounds are one. This proves the stated gap and
\eqref{eq:asymptotic}.
It also gives, uniformly in all admissible parameters,
\begin{equation}\label{eq:uniform-scale}
 c(e,n)=\Theta\!\left(\max\!\left\{1,\frac{n}{e^2+1}\right\}\right).
\end{equation}
For example, absolute comparison constants $1/2$ and $2$ follow
directly from \eqref{eq:main}.

The bounds recover $c(1,n)=\lceil n/2\rceil$, also a consequence of
the Hermitian higher-rank numerical-range formula~\cite{CKZ}.
They give $c(e,n)=1$ whenever $n\le e^2+1$, and exact values at
$n=q(e^2+1)$ and $n=q(e^2+1)-1$ for every admissible positive $q$.
For two errors, this yields $c(2,9)=c(2,10)=2$ and
$c(2,14)=c(2,15)=3$. At $(e,n)=(2,6)$, the present argument leaves
$1\le c(2,6)\le2$.

The corresponding integer question for arbitrary operator systems is
whether every $D$-dimensional system has a scalar compression of rank
$\lceil n/D\rceil$. This is a form of Weaver's conjecture on the lower
quantum Tur\'an numbers~\cite[following Theorem~2.11]{Weaver}.
For $D=3$, the stronger threshold $n\ge3k-2$ is known~\cite[Lemma~2.3]{AK}.
Theorem~\ref{thm:compression} gives $n\ge Dk-1$ for every $D$ but does
not reach the conjectured threshold $n\ge D(k-1)+1$ in general.
The remaining rounding problem for error spaces may also benefit from
their factorization as products $A_i^*A_j$, which was not used beyond
the dimension bound in Corollary~\ref{cor:lower}.

Finally, normalization changes the extremal problem. If a chosen error
basis satisfies $\sum_i A_i^*A_i=I_n$, then the identity already lies
in the product span and $\dim\V_{\E}\le e^2$. The scalar-compression
theorem still applies with that smaller dimension. The explicit models
in Section~\ref{sec:models}, however, use a common zero block and need
not satisfy this normalization. Equation~\eqref{eq:main} concerns the
arbitrary linear error spaces in \eqref{eq:minimax}.


\section*{Verification and research disclosure}

The manuscript has received sequential internal analytical review within
the same AI-assisted research workflow. This is not independent expert
acceptance or formal proof-assistant certification. The complete arguments,
including the metric convention, phase-summed Hessian, singular Gram
matrices, channel normalization and integer thresholds, were checked
against their stated hypotheses. During integration, an auxiliary claim
about a literal left kernel was restricted to $n>(m+1)k$; the main
quantitative theorem already applies at $n\ge(m+1)k-1$ and is unchanged.

The accompanying source package contains executable finite checks. The
inherited checks reconstruct nine exact second-variation identities, reject
nine deliberately invalid orthogonality controls, and exercise the explicit
block models. The new checks cover boundary Hessians, sharp traceless-stack
bounds, phase sums, error-basis conversion, singular and nonsingular
recovery, a rational dephasing example, complex syndrome conventions,
and excluded normalization and dimension cases. Four floating-point
search runs illustrate both descent branches and the encoding--recovery
pipeline. Their small residuals are diagnostics, not certified interval
bounds or proofs of universal correctness. The universal statements rest
on the analytic arguments above.

The derivation, primary-source comparison, computations and manuscript
preparation used AI assistance under human direction. Independent expert
review, comprehensive novelty assessment, and a certified finite-precision
implementation remain outstanding. Version 2 retains the earlier exact
asymptotic result and adds the quantitative construction; it does not
claim a new resolution of the remaining integer problem.

\begin{thebibliography}{99}
\raggedright
\bibitem{KW05}
O.~Kr\"uger and R.~F.~Werner (eds.),
\emph{Some Open Problems in Quantum Information Theory} (2005),
Problem~14, \href{https://arxiv.org/abs/quant-ph/0504166}{arXiv:quant-ph/0504166}.

\bibitem{QIQC}
QIQCOP Zoo, \emph{Tough error models}, problem record
\href{https://qiqc-op.com/problem/op_b1b41f737f9b4aeb/}{op\_b1b41f737f9b4aeb},
version dated 4 September 2026; accessed 16 November 2026.

\bibitem{KL97}
E.~Knill and R.~Laflamme, Theory of quantum error-correcting codes,
\emph{Physical Review A} \textbf{55}, 900--911 (1997),
\href{https://doi.org/10.1103/PhysRevA.55.900}{doi:10.1103/PhysRevA.55.900}.

\bibitem{KLV00}
E.~Knill, R.~Laflamme, and L.~Viola,
Theory of quantum error correction for general noise,
\emph{Physical Review Letters} \textbf{84}, 2525--2528 (2000),
\href{https://doi.org/10.1103/PhysRevLett.84.2525}{doi:10.1103/PhysRevLett.84.2525};
\href{https://arxiv.org/abs/quant-ph/9908066}{arXiv:quant-ph/9908066}.

\bibitem{Reimpell}
M.~Reimpell, \emph{Quantum Information and Convex Optimization},
doctoral dissertation, Technische Universit\"at Braunschweig (2008),
Section~3.3.4,
\href{https://www.osti.gov/etdeweb/servlets/purl/21161752}{archived dissertation}.

\bibitem{Weaver}
N.~Weaver, The ``quantum'' Tur\'an problem for operator systems,
\emph{Pacific Journal of Mathematics} \textbf{301}, 335--349 (2019),
\href{https://doi.org/10.2140/pjm.2019.301.335}{doi:10.2140/pjm.2019.301.335};
\href{https://arxiv.org/abs/1802.07394}{arXiv:1802.07394}.

\bibitem{AK}
A.~Allen and A.~Kornell, The quantum Ramsey numbers $QR(2,k)$,
\emph{Linear and Multilinear Algebra} \textbf{73}, 3959--3964 (2025),
\href{https://doi.org/10.1080/03081087.2025.2578529}{doi:10.1080/03081087.2025.2578529};
\href{https://arxiv.org/abs/2507.04128v2}{arXiv:2507.04128v2}.

\bibitem{LiTam}
C.-K.~Li and B.-S.~Tam, A note on extreme correlation matrices,
\emph{SIAM Journal on Matrix Analysis and Applications} \textbf{15},
903--908 (1994),
\href{https://doi.org/10.1137/S0895479892240683}{doi:10.1137/S0895479892240683}.

\bibitem{Gupta}
A.~Gupta, B.~Sun, X.~He, and B.~Zeng,
\emph{Variance Geometry of Exact Pauli-Detecting Codes:
Continuous Landscapes Beyond Stabilizers} (2026),
\href{https://arxiv.org/abs/2604.21800v1}{arXiv:2604.21800v1}.

\bibitem{Sun}
B.~Sun, N.~Cao, Y.~T.~Poon, and B.~Zeng,
\emph{Geometry of Knill--Laflamme Coefficients for Pauli Error Detection}
(2026), \href{https://arxiv.org/abs/2610.01992v1}{arXiv:2610.01992v1}.

\bibitem{CKZ}
M.-D.~Choi, D.~W.~Kribs, and K.~\.{Z}yczkowski,
Higher-rank numerical ranges and compression problems,
\emph{Linear Algebra and its Applications} \textbf{418}, 828--839 (2006),
\href{https://doi.org/10.1016/j.laa.2006.03.019}{doi:10.1016/j.laa.2006.03.019};
\href{https://arxiv.org/abs/math/0511278}{arXiv:math/0511278}.

\bibitem{NgMandayam}
H.~K.~Ng and P.~Mandayam,
A simple approach to approximate quantum error correction based on the
transpose channel, \emph{Physical Review A} \textbf{81}, 062342 (2010),
\href{https://arxiv.org/abs/0909.0931v2}{arXiv:0909.0931v2}.

\bibitem{BenyOreshkov}
C.~B\'eny and O.~Oreshkov,
General conditions for approximate quantum error correction and near-optimal
recovery channels, \emph{Physical Review Letters} \textbf{104}, 120501 (2010),
\href{https://arxiv.org/abs/0907.5391}{arXiv:0907.5391}.

\bibitem{KSW}
D.~Kretschmann, D.~Schlingemann, and R.~F.~Werner,
A continuity theorem for Stinespring's dilation,
\href{https://arxiv.org/abs/0710.2495}{arXiv:0710.2495} (2007).

\end{thebibliography}
\end{document}
